% ============================================================
% f2014_anhui_science_q03_flowchart.tex
% 2014 安徽卷（理科）第 3 题：程序框图
% 来源 PDF：普通高考/2014/2014安徽理.pdf
% 原页：第 1 页（A3 横排）左栏，印面 857，第 3 题题干下方
% 结构：开始 → 赋初值 x=1、y=1 → z=x+y → 判断 z≤50；
%       是：x=y → y=z → 返回 z=x+y 之前；否：输出 z → 结束。
% ============================================================
\begin{tikzpicture}[>=Stealth, line width=0.7pt, font=\small,
  x=1cm, y=1cm, line cap=round, line join=round]

  % 开始
  \draw[rounded corners=6pt] (-0.95,0.42) rectangle (0.95,-0.42);
  \node at (0,0) {开始};
  \draw[->] (0,-0.42) -- (0,-0.85);

  % 赋初值
  \draw (-1.55,-0.85) rectangle (1.55,-1.70);
  \node at (0,-1.28) {$x=1$，$y=1$};
  \draw[->] (0,-1.70) -- (0,-2.10);

  % z=x+y
  \draw (-1.30,-2.10) rectangle (1.30,-2.95);
  \node at (0,-2.52) {$z=x+y$};
  \draw[->] (0,-2.95) -- (0,-3.40);

  % 判断框
  \draw (0,-3.40) -- (1.80,-4.30) -- (0,-5.20) -- (-1.80,-4.30) -- cycle;
  \node at (0,-4.30) {$z\leqslant 50$};

  % 是：向下进入循环体
  \draw[->] (0,-5.20) -- (0,-5.60);
  \node at (0.58,-5.42) {是};
  \draw (-1.20,-5.60) rectangle (1.20,-6.45);
  \node at (0,-6.02) {$x=y$};

  \draw[->] (0,-6.45) -- (0,-6.85);
  \draw (-1.20,-6.85) rectangle (1.20,-7.70);
  \node at (0,-7.27) {$y=z$};

  % 回边：向下、向左、向上，再水平汇入 z=x+y 上方
  \draw (-1.20,-7.27) -- (-2.95,-7.27) -- (-2.95,-1.90) -- (0,-1.90);
  \draw[->] (0,-1.90) -- (0,-2.10);

  % 否：向右输出
  \draw[->] (1.80,-4.30) -- (2.98,-4.30);
  \node at (2.35,-3.98) {否};
  \draw (3.15,-3.88) -- (5.35,-3.88) -- (5.05,-4.72) -- (2.85,-4.72) -- cycle;
  \node at (4.10,-4.30) {输出 $z$};
  \draw[->] (4.10,-4.72) -- (4.10,-5.10);

  % 结束
  \draw[rounded corners=6pt] (3.15,-5.10) rectangle (5.05,-5.94);
  \node at (4.10,-5.52) {结束};

\end{tikzpicture}
